\documentclass[a4paper,11pt]{article}
\usepackage[utf8]{inputenc}
\usepackage[T1]{fontenc}
\usepackage{geometry}
\usepackage{enumitem}
\usepackage{titlesec}
\usepackage[table]{xcolor}
\usepackage{fancyhdr}
\usepackage{tikz}
\usetikzlibrary{shapes,arrows,positioning}
\usepackage{tabularx}
\usepackage{listings}
\usepackage{graphicx}
\usepackage{lastpage}
\usepackage{hyperref}
\usepackage{float}

% Page Setup
\geometry{top=1in, bottom=1in, left=0.8in, right=0.8in}
\pagestyle{fancy}
\fancyhf{}
\lhead{\small Project 83: Agentic AI Compiler Assistant}
\rhead{\small Comprehensive Documentation}
\cfoot{\thepage\ of \pageref{LastPage}}

% Formatting
\titleformat{\section}{\Large\bfseries\color{blue!70!black}}{\thesection.}{1em}{}[\titlerule]
\titleformat{\subsection}{\large\bfseries\color{blue!50!black}}{\thesubsection}{1em}{}

% Code Listings Configuration
\lstset{
    basicstyle=\ttfamily\small,
    breaklines=true,
    frame=single,
    backgroundcolor=\color{gray!10},
    keywordstyle=\color{blue},
    stringstyle=\color{green!50!black},
    commentstyle=\color{gray},
    showstringspaces=false,
    tabsize=4,
    captionpos=b
}

% Custom commands
\newcommand{\statusgood}[1]{\colorbox{green!20}{\textbf{#1}}}
\newcommand{\statuspending}[1]{\colorbox{yellow!20}{\textbf{#1}}}
\newcommand{\statusissue}[1]{\colorbox{red!20}{\textbf{#1}}}

\begin{document}

\begin{center}
    {\Huge \textbf{Comprehensive Project Documentation}} \\
    \vspace{5mm}
    {\LARGE \textbf{Project 83: Conversational Agentic AI}} \\
    {\LARGE \textbf{Compiler Assistant for Mini-Python}} \\
    \vspace{3mm}
    {\large \textit{A Novel Approach to Intelligent Compiler Design}} \\
    \vspace{3mm}
    \textbf{Status:} \statusgood{PHASE 4 COMPLETED} \\
    \textbf{Current Week:} Week 7 | \textbf{Implementation:} Parser Complete
\end{center}

\vspace{5mm}

\tableofcontents
\clearpage

\section{Problem Statement}

\subsection{1.1 Background and Motivation}

Traditional compilers, while powerful and efficient, present significant challenges to both novice programmers and students learning compiler design. When compilation errors occur, conventional compilers typically provide:

\begin{itemize}[leftmargin=*]
    \item \textbf{Cryptic error messages:} Technical jargon without context (e.g., "syntax error: expected ':'")
    \item \textbf{Limited guidance:} No explanation of \textit{why} the code is incorrect
    \item \textbf{No suggestions:} Users must independently determine how to fix errors
    \item \textbf{Single error focus:} Compilation halts at first major error, hiding subsequent issues
    \item \textbf{Steep learning curve:} Students struggle to connect error messages to underlying grammar rules
\end{itemize}

\subsection{1.2 Current Limitations in Compiler Technology}

\textbf{For Traditional Compilers (GCC, Python's CPython):}
\begin{itemize}
    \item Provide error codes and line numbers only
    \item No natural language explanations
    \item No interactive debugging assistance
    \item Limited educational value for learning compiler theory
\end{itemize}

\textbf{For Code Editors and IDEs (VS Code, PyCharm):}
\begin{itemize}
    \item Syntax highlighting helps but doesn't explain errors
    \item Some provide autocomplete, but not grammar-based reasoning
    \item Require manual consultation of documentation
    \item Not integrated with formal language specifications (EBNF)
\end{itemize}

\textbf{For AI Code Assistants (GitHub Copilot, ChatGPT):}
\begin{itemize}
    \item General-purpose tools, not compiler-specific
    \item Lack integration with compilation pipeline
    \item Cannot access token streams or AST structures
    \item No grounding in formal grammar rules
    \item Suggestions may not align with language specifications
\end{itemize}

\subsection{1.3 The Gap This Project Addresses}

There is a \textbf{critical gap} between:
\begin{enumerate}
    \item \textbf{Compiler error detection} (what current compilers do well), and
    \item \textbf{Intelligent error explanation and guidance} (what users actually need)
\end{enumerate}

\textbf{Specific Problems Solved:}

\begin{table}[H]
\centering
\begin{tabularx}{\textwidth}{|l|X|X|}
\hline
\rowcolor{blue!10}
\textbf{Problem} & \textbf{Current State} & \textbf{Our Solution} \\
\hline
Understanding Errors & "Missing colon" & "Function definitions in Mini-Python require a colon (:) before the indented block, per grammar rule..." \\
\hline
Learning Grammar & Manual EBNF study & AI references specific grammar rules in plain language \\
\hline
Security Awareness & No warnings & Proactive flagging of insecure patterns (eval, exec) \\
\hline
Debugging Process & Trial and error & Interactive conversation with AI for iterative refinement \\
\hline
Educational Value & Minimal & Teaches compiler theory through hands-on interaction \\
\hline
\end{tabularx}
\caption{Problems Addressed by This Project}
\end{table}

\subsection{1.4 Target Audience and Use Cases}

\textbf{Primary Users:}
\begin{itemize}
    \item \textbf{Computer Science Students:} Learning compiler design and formal languages
    \item \textbf{Novice Programmers:} Understanding syntax and language rules
    \item \textbf{Educators:} Teaching compiler theory with visual, interactive tools
    \item \textbf{Researchers:} Exploring AI applications in program analysis
\end{itemize}

\textbf{Use Cases:}
\begin{enumerate}
    \item Student debugging Mini-Python code with immediate AI feedback
    \item Professor demonstrating compilation phases (Lexer → Parser → AST) in real-time
    \item Developer prototyping code with security audit suggestions
    \item Self-learner studying how formal grammars map to real code
\end{enumerate}

\subsection{1.5 Project Objectives}

\textbf{Primary Objective:} Develop an intelligent compiler for Mini-Python that combines traditional compilation phases with AI-powered assistance to provide educational, interactive, and secure code analysis.

\textbf{Specific Goals:}
\begin{enumerate}
    \item Implement a complete compiler front-end (Lexer, Parser, Semantic Analyzer) from scratch
    \item Integrate Google Gemini AI to analyze errors and suggest corrections
    \item Provide natural language explanations referencing EBNF grammar rules
    \item Create conversational interface for iterative debugging
    \item Build security auditing layer to identify dangerous code patterns
    \item Achieve >95\% test coverage for reliability
    \item Develop both CLI and Web UI for accessibility
\end{enumerate}

\clearpage

\section{Novelty and Innovation}

\subsection{2.1 Core Innovation}

This project introduces a \textbf{paradigm shift} in compiler design by integrating \textbf{Agentic AI} into the traditional compilation pipeline. Unlike conventional compilers that simply report errors, this system employs Google Gemini's advanced language model to provide \textbf{intelligent, conversational assistance} throughout the compilation process.

\subsection{2.2 Novel Contributions}

\begin{enumerate}[leftmargin=*]
    \item \textbf{AI-Augmented Error Recovery:}
    \begin{itemize}
        \item Traditional compilers terminate on first major error
        \item Our system uses AI to analyze error context and suggest corrections
        \item Gemini API provides natural language explanations of syntax violations
        \item References EBNF grammar rules to justify corrections
    \end{itemize}
    
    \item \textbf{Conversational Debugging Interface:}
    \begin{itemize}
        \item Interactive dialogue between user and AI agent
        \item Multi-turn conversation for iterative code refinement
        \item Context-aware suggestions based on entire program structure
        \item Real-time feedback loop for learning compiler concepts
    \end{itemize}
    
    \item \textbf{Proactive Security Analysis:}
    \begin{itemize}
        \item AI identifies insecure coding patterns during compilation
        \item Flags dangerous constructs (e.g., \texttt{eval()}, \texttt{exec()})
        \item Suggests secure alternatives with explanations
        \item Security auditing integrated into semantic analysis phase
    \end{itemize}
    
    \item \textbf{Educational Value:}
    \begin{itemize}
        \item System explains WHY code is incorrect, not just WHAT is wrong
        \item References formal grammar rules in plain language
        \item Helps students understand compiler theory through interaction
        \item Bridges theory and practice in compiler design education
    \end{itemize}
\end{enumerate}

\subsection{2.3 Technical Innovations}

\begin{table}[H]
\centering
\begin{tabularx}{\textwidth}{|l|X|}
\hline
\rowcolor{blue!10}
\textbf{Innovation} & \textbf{Description} \\
\hline
Hybrid Architecture & Manual compiler front-end + AI-powered assistance layer \\
\hline
Grammar-Aware AI & Gemini model injected with EBNF grammar as system instruction \\
\hline
Dynamic Model Selection & Auto-detection of available Gemini models via API \\
\hline
AST-to-Prompt Translation & Converts Abstract Syntax Tree to natural language for AI \\
\hline
Error Context Enrichment & Provides AI with token stream, line numbers, and scope info \\
\hline
\end{tabularx}
\caption{Technical Innovation Summary}
\end{table}

\subsection{2.4 Differentiators from Existing Systems}

\textbf{Compared to Traditional Compilers (GCC, Clang):}
\begin{itemize}
    \item Provides explanations, not just error codes
    \item Suggests fixes in natural language
    \item Interactive refinement capability
\end{itemize}

\textbf{Compared to AI Code Assistants (GitHub Copilot, ChatGPT):}
\begin{itemize}
    \item Deeply integrated with compilation pipeline
    \item Uses formal grammar as ground truth
    \item Compiler-specific reasoning (not general coding help)
    \item Educational focus on compiler theory
\end{itemize}

\clearpage

\section{Proposed System Architecture}

\subsection{3.1 High-Level Architecture Overview}

The system follows a \textbf{three-layer architecture}:
\begin{enumerate}
    \item \textbf{Compiler Front-End Layer:} Manual implementation of Lexer, Parser, and Semantic Analyzer
    \item \textbf{Agentic AI Layer:} Gemini API integration with specialized prompts
    \item \textbf{User Interface Layer:} CLI and Web UI (Streamlit) for interaction
\end{enumerate}

\subsection{3.2 Architecture Diagram}

\textbf{Note:} The complete architecture diagram is shown below. For PlantUML source code to regenerate this diagram, see subsection 3.3.

\begin{figure}[H]
\centering
% TODO: Add your architecture image here
% Uncomment and replace with your image filename:
% \includegraphics[width=0.95\textwidth]{architecture_diagram.png}
\fbox{\parbox{0.95\textwidth}{\centering\vspace{3cm}\Large\textbf{[INSERT ARCHITECTURE DIAGRAM IMAGE HERE]}\\\vspace{0.5cm}\normalsize\textit{Place your compiled PlantUML architecture diagram in this location}\vspace{3cm}}}
\caption{Agentic AI Compiler System Architecture}
\label{fig:architecture}
\end{figure}

\subsection{3.3 PlantUML Architecture Code}

The following PlantUML code generates the architecture diagram:}

\begin{lstlisting}[language=bash, caption=PlantUML Architecture Code]
@startuml
!define RECTANGLE class

skinparam componentStyle rectangle
skinparam linetype ortho

package "User Interface Layer" {
    [Web UI (Streamlit)] as WebUI
    [Command Line Interface] as CLI
    note right of WebUI
        - Code Editor Panel
        - AI Chatbot Panel
        - AST Visualization
    end note
}

package "Compiler Front-End Layer" {
    [Lexer] as Lexer
    [Parser] as Parser
    [Semantic Analyzer] as Semantic
    [AST Builder] as AST
    [Error Diagnostics] as Diagnostics
    
    note right of Lexer
        Manual Implementation
        - Regex-based tokenization
        - Indentation handling
        - Token stream generation
    end note
}

package "Agentic AI Layer" {
    [Orchestrator] as Orchestrator
    [Gemini Client] as Gemini
    [Prompt Engineer] as Prompts
    [Tool Definitions] as Tools
    
    note right of Orchestrator
        The "Brain"
        - Manages feedback loops
        - Routes errors to AI
        - Formats responses
    end note
}

package "Data \& Configuration" {
    database "EBNF Grammar" as Grammar
    database "System Prompts" as SysPrompts
    database ".env (API Keys)" as Env
    database "Test Cases" as Tests
}

' User Interaction Flow
WebUI --> Orchestrator : Submit Code
CLI --> Orchestrator : Submit Code

' Compilation Pipeline
Orchestrator --> Lexer : Source Code
Lexer --> Parser : Token Stream
Parser --> AST : Build AST
AST --> Semantic : Analyze
Semantic --> Diagnostics : Errors/Warnings

' AI Integration Flow
Diagnostics --> Orchestrator : Error Context
Orchestrator --> Gemini : Enhanced Prompt
Gemini --> Orchestrator : AI Response
Orchestrator --> WebUI : Display Suggestions
Orchestrator --> CLI : Display Suggestions

' Data Dependencies
Grammar --> Prompts : Grammar Rules
SysPrompts --> Gemini : Context Injection
Env --> Gemini : API Authentication
Tools --> Gemini : Function Calling

' Testing
Tests --> Lexer : Validate
Tests --> Gemini : Verify AI

@enduml
\end{lstlisting}

\subsection{3.4 Component Interaction Sequence}

\textbf{Normal Compilation Flow (No Errors):}
\begin{enumerate}
    \item User submits Mini-Python code via UI
    \item Orchestrator passes code to Lexer
    \item Lexer tokenizes and returns token stream
    \item Parser builds AST from tokens
    \item Semantic Analyzer validates AST
    \item Success message returned to user
\end{enumerate}

\textbf{Error Handling Flow (With AI Assistance):}
\begin{enumerate}
    \item Compilation stops at first error (Lexer/Parser/Semantic)
    \item Diagnostics module creates detailed error context:
    \begin{itemize}
        \item Error type and message
        \item Line and column numbers
        \item Surrounding code snippet
        \item Current token stream state
    \end{itemize}
    \item Orchestrator enriches error with EBNF grammar rules
    \item Gemini Client sends prompt to AI model
    \item AI analyzes error against grammar and returns:
    \begin{itemize}
        \item Explanation of what went wrong
        \item Suggested correction
        \item Reference to violated grammar rule
    \end{itemize}
    \item Orchestrator formats response for UI
    \item User can refine code and resubmit
\end{enumerate}

\clearpage

\section{Prerequisites and System Requirements}

\subsection{4.1 Hardware Requirements}

\begin{table}[H]
\centering
\begin{tabularx}{\textwidth}{|l|X|X|}
\hline
\rowcolor{blue!10}
\textbf{Component} & \textbf{Minimum} & \textbf{Recommended} \\
\hline
Processor & Intel Core i3 / AMD Ryzen 3 & Intel Core i5 / AMD Ryzen 5 \\
\hline
RAM & 4 GB & 8 GB or higher \\
\hline
Storage & 500 MB free space & 1 GB free space \\
\hline
Internet & Stable connection (for Gemini API) & High-speed broadband \\
\hline
\end{tabularx}
\caption{Hardware Requirements}
\end{table}

\subsection{4.2 Software Prerequisites}

\textbf{Core Requirements:}
\begin{itemize}[leftmargin=*]
    \item \textbf{Python 3.10+} (Tested on Python 3.12)
    \begin{itemize}
        \item Download: \url{https://www.python.org/downloads/}
        \item Verify installation: \texttt{python --version}
    \end{itemize}
    
    \item \textbf{pip} (Python Package Manager)
    \begin{itemize}
        \item Usually included with Python
        \item Verify: \texttt{pip --version}
    \end{itemize}
    
    \item \textbf{Google Gemini API Key}
    \begin{itemize}
        \item Sign up: \url{https://makersuite.google.com/app/apikey}
        \item Free tier available with usage limits
    \end{itemize}
\end{itemize}

\textbf{Python Dependencies:}
\begin{lstlisting}[language=bash, caption=requirements.txt]
google-generativeai>=0.3.0
python-dotenv>=1.0.0
pytest>=7.4.0
\end{lstlisting}

\textbf{Installation Command:}
\begin{lstlisting}[language=bash]
pip install -r requirements.txt
\end{lstlisting}

\subsection{4.3 Development Tools (Optional)}

\begin{itemize}[leftmargin=*]
    \item \textbf{VS Code} / \textbf{PyCharm} - Recommended IDEs
    \item \textbf{Git} - Version control
    \item \textbf{Graphviz} - For AST visualization (planned feature)
    \item \textbf{LaTeX} - For report generation
\end{itemize}

\subsection{4.4 Knowledge Prerequisites}

\textbf{For Users:}
\begin{itemize}
    \item Basic understanding of Python syntax
    \item Familiarity with command line / terminal
\end{itemize}

\textbf{For Developers:}
\begin{itemize}
    \item Compiler theory basics (Lexer, Parser, AST)
    \item Regular expressions (for Lexer implementation)
    \item Understanding of EBNF grammar notation
    \item API integration concepts
    \item Python object-oriented programming
\end{itemize}

\clearpage

\section{Tools and Technologies Used}

\subsection{5.1 Implementation Tools}

\begin{table}[H]
\centering
\begin{tabularx}{\textwidth}{|l|l|X|}
\hline
\rowcolor{blue!10}
\textbf{Tool/Technology} & \textbf{Version} & \textbf{Purpose} \\
\hline
Python & 3.10+ & Core programming language \\
\hline
google-generativeai & 0.3.0+ & Gemini API SDK \\
\hline
python-dotenv & 1.0.0+ & Environment variable management \\
\hline
pytest & 7.4.0+ & Unit testing framework \\
\hline
Regex (re module) & Built-in & Token pattern matching \\
\hline
Enum & Built-in & Token type definitions \\
\hline
Streamlit & 1.28+ (Planned) & Web UI framework \\
\hline
\end{tabularx}
\caption{Primary Tools and Libraries}
\end{table}

\subsection{5.2 Google Gemini AI}

\textbf{Why Gemini?}
\begin{itemize}[leftmargin=*]
    \item \textbf{128k Token Context Window:} Allows feeding entire EBNF grammar + error context
    \item \textbf{System Instructions:} Persistent persona as "Compiler Assistant"
    \item \textbf{Function Calling:} Planned feature for AST inspection
    \item \textbf{Fast Response Time:} Critical for interactive debugging
    \item \textbf{Free Tier Available:} Suitable for educational projects
\end{itemize}

\textbf{Model Selection:}
\begin{lstlisting}[language=Python, caption=Dynamic Model Discovery]
import google.generativeai as genai

genai.configure(api_key=os.getenv('GEMINI_API_KEY'))

# Auto-detect available models
available_models = genai.list_models()
for model in available_models:
    if 'generateContent' in model.supported_generation_methods:
        print(f"Using: {model.name}")
        break
\end{lstlisting}

\subsection{5.3 Development Environment}

\textbf{Directory Structure:}
\begin{lstlisting}[language=bash]
CD-Project/
├── src/
│   ├── compiler/
│   │   ├── tokens.py          # Token definitions
│   │   └── lexer.py           # Lexer implementation
│   ├── orchestrator/
│   │   └── gemini_client.py   # AI integration
│   └── tests/
│       └── ai_test_cases.py   # AI validation
├── grammar/
│   └── mini_python.ebnf       # Formal grammar
├── docs/
│   ├── SRS.md                 # Requirements
│   ├── prompts.md             # AI prompts
│   └── tech_stack.md          # Technology docs
├── reports/
│   └── week_05_report.tex     # Progress reports
├── .env                       # API keys (GITIGNORED)
├── requirements.txt           # Dependencies
└── manual_test.py             # Integration tests
\end{lstlisting}

\subsection{5.4 Version Control and Collaboration}

\begin{itemize}[leftmargin=*]
    \item \textbf{Git} for version control
    \item \textbf{.gitignore} configured to exclude:
    \begin{itemize}
        \item \texttt{.env} (API keys)
        \item \texttt{\_\_pycache\_\_/} (Python bytecode)
        \item \texttt{venv/} (Virtual environment)
    \end{itemize}
\end{itemize}

\clearpage

\section{Preliminary Implementation Completed}

\subsection{6.1 Phase 1: Project Setup (Weeks 1-2)}

\textbf{Completed Deliverables:}
\begin{itemize}[leftmargin=*]
    \item \statusgood{COMPLETE} Project structure established
    \item \statusgood{COMPLETE} EBNF grammar for Mini-Python defined
    \item \statusgood{COMPLETE} Gemini API integration tested
    \item \statusgood{COMPLETE} Dynamic model selection implemented
    \item \statusgood{COMPLETE} Environment configuration (.env setup)
\end{itemize}

\textbf{Key Achievement:} Successfully authenticated with Gemini API and verified AI reasoning capabilities.

\subsection{6.2 Phase 2: AI Foundation (Weeks 3-4)}

\textbf{Completed Deliverables:}
\begin{itemize}[leftmargin=*]
    \item \statusgood{COMPLETE} \texttt{gemini\_client.py} implementation
    \item \statusgood{COMPLETE} System prompt engineering
    \item \statusgood{COMPLETE} AI test cases for syntax error detection
    \item \statusgood{COMPLETE} Security auditing test cases
    \item \statusgood{COMPLETE} SRS documentation
    \item \statusgood{COMPLETE} Architecture design documents
\end{itemize}

\textbf{AI Capabilities Verified:}
\begin{enumerate}
    \item Detect missing colons in function definitions
    \item Flag insecure \texttt{eval()} usage
    \item Provide grammar-based explanations
    \item Suggest corrections with reasoning
\end{enumerate}

\subsection{6.3 Phase 3: Lexer Implementation (Week 5)}

\textbf{Status:} \statusgood{COMPLETED}

\subsubsection{5.3.1 Token Type System}

\begin{lstlisting}[language=Python, caption=Token Definition]
from enum import Enum, auto

class TokenType(Enum):
    # Keywords
    DEF = auto()
    IF = auto()
    ELIF = auto()
    ELSE = auto()
    WHILE = auto()
    FOR = auto()
    RETURN = auto()
    IN = auto()
    AND = auto()
    OR = auto()
    NOT = auto()
    
    # Literals
    INTEGER = auto()
    STRING = auto()
    TRUE = auto()
    FALSE = auto()
    
    # Operators
    PLUS = auto()
    MINUS = auto()
    MULTIPLY = auto()
    DIVIDE = auto()
    MODULO = auto()
    EQUALS = auto()
    EQ = auto()      # ==
    NEQ = auto()     # !=
    LT = auto()      # <
    GT = auto()      # >
    LTE = auto()     # <=
    GTE = auto()     # >=
    
    # Delimiters
    LPAREN = auto()
    RPAREN = auto()
    COLON = auto()
    COMMA = auto()
    
    # Structural
    INDENT = auto()
    DEDENT = auto()
    NEWLINE = auto()
    
    # Special
    IDENTIFIER = auto()
    EOF = auto()
    
class Token:
    def __init__(self, type: TokenType, value: str, 
                 line: int, column: int):
        self.type = type
        self.value = value
        self.line = line
        self.column = column
\end{lstlisting}

\subsubsection{6.3.2 Lexer Core Implementation}

\textbf{Key Features:}
\begin{itemize}[leftmargin=*]
    \item \textbf{Regex-based pattern matching} for efficient tokenization
    \item \textbf{Single-pass scanning} with position tracking
    \item \textbf{Line and column number recording} for error reporting
    \item \textbf{Indentation stack management} for Python-style blocks
    \item \textbf{Comprehensive error handling} with descriptive messages
\end{itemize}

\textbf{Implementation Statistics:}
\begin{itemize}
    \item \textbf{Lines of Code:} 320 lines (lexer.py)
    \item \textbf{Token Types:} 36 distinct types
    \item \textbf{Regex Patterns:} 25 patterns
    \item \textbf{Test Cases:} 45 tests
    \item \textbf{Code Coverage:} 98\%
\end{itemize}

\subsubsection{6.3.3 Indentation Handling Algorithm}

\textbf{Challenge:} Python uses significant whitespace (indentation) to define code blocks, unlike braces in C/Java.

\textbf{Solution:} Stack-based indentation tracking
\begin{enumerate}
    \item Maintain \texttt{indent\_stack} initialized with [0]
    \item For each line:
    \begin{itemize}
        \item Count leading spaces
        \item Compare with \texttt{indent\_stack[-1]}
        \item If increased: emit INDENT, push new level
        \item If decreased: emit DEDENT(s), pop until match
        \item If unchanged: no tokens
    \end{itemize}
    \item At EOF: emit DEDENT for each remaining stack level
\end{enumerate}

\begin{lstlisting}[language=Python, caption=Indentation Handler]
def _handle_indentation(self, line: str) -> List[Token]:
    if not line.strip():
        return []
    
    spaces = len(line) - len(line.lstrip())
    current_indent = self.indent_stack[-1]
    tokens = []
    
    if spaces > current_indent:
        self.indent_stack.append(spaces)
        tokens.append(Token(TokenType.INDENT, '', 
                          self.current_line, 0))
    
    elif spaces < current_indent:
        while self.indent_stack[-1] > spaces:
            self.indent_stack.pop()
            tokens.append(Token(TokenType.DEDENT, '', 
                              self.current_line, 0))
        
        if self.indent_stack[-1] != spaces:
            raise IndentationError(
                f"Unindent does not match any outer level "
                f"at line {self.current_line}")
    
    return tokens
\end{lstlisting}

\subsection{6.4 Testing and Validation}

\textbf{Test Categories:}
\begin{table}[H]
\centering
\begin{tabularx}{\textwidth}{|l|c|X|}
\hline
\rowcolor{blue!10}
\textbf{Category} & \textbf{Tests} & \textbf{Coverage} \\
\hline
Basic Tokenization & 10 & Keywords, literals, operators \\
\hline
Identifiers & 5 & Valid/invalid variable names \\
\hline
Indentation & 8 & Nested blocks, mixed indent errors \\
\hline
Error Handling & 12 & Invalid tokens, unterminated strings \\
\hline
Complex Code & 10 & Full functions, loops, conditionals \\
\hline
\textbf{Total} & \textbf{45} & \textbf{98\% Line Coverage} \\
\hline
\end{tabularx}
\caption{Test Suite Summary}
\end{table}

\textbf{Performance Metrics:}
\begin{itemize}
    \item \textbf{Tokenization Speed:} 1,250 lines/second (Target: >1,000)
    \item \textbf{Memory Usage:} ~50 bytes per token
    \item \textbf{Largest File Tested:} 500 lines (Factorial examples)
\end{itemize}

\subsection{6.5 Parser Implementation (Weeks 6-7)}

\textbf{Status:} \statusgood{COMPLETED}

\subsubsection{6.5.1 AST Node Hierarchy}

Implemented comprehensive AST node system with 16 node types covering all Mini-Python constructs:

\textbf{Node Categories:}
\begin{itemize}[leftmargin=*]
    \item \textbf{Base Classes:} ASTNode, Expression, Statement
    \item \textbf{Expressions:} BinaryOp, UnaryOp, Identifier, Literal, FunctionCall
    \item \textbf{Statements:} Assignment, IfStatement, WhileStatement, FunctionDef, ReturnStatement, PrintStatement
    \item \textbf{Structural:} Block, Program
\end{itemize}

\subsubsection{6.5.2 Recursive Descent Parser}

Implemented full parser with:
\begin{itemize}[leftmargin=*]
    \item \textbf{Expression Parsing:} Correct operator precedence (*, / before +, -)
    \item \textbf{Statement Parsing:} All statement types including control flow
    \item \textbf{Block Handling:} Proper INDENT/DEDENT token consumption
    \item \textbf{Error Reporting:} Position-aware error messages
\end{itemize}

\textbf{Parser Statistics:}
\begin{itemize}[leftmargin=*]
    \item \textbf{Lines of Code:} 820 lines (parser.py + ast\_nodes.py)
    \item \textbf{Test Cases:} 33 tests (27 parser + 6 integration)
    \item \textbf{Test Pass Rate:} 100\%
    \item \textbf{Code Coverage:} 96\%
\end{itemize}

\subsection{6.6 Current Status and Next Steps}

\begin{table}[H]
\centering
\begin{tabularx}{\textwidth}{|l|X|}
\hline
\rowcolor{blue!10}
\textbf{Component} & \textbf{Status} \\
\hline
Lexer & \statusgood{COMPLETE} - Week 5 \\
\hline
Parser & \statusgood{COMPLETE} - Weeks 6-7 \\
\hline
Semantic Analysis & \statuspending{PENDING} - Week 8 \\
\hline
AI Orchestration & \statuspending{PENDING} - Week 9 \\
\hline
Web UI & \statuspending{PENDING} - Weeks 10-11 \\
\hline
\end{tabularx}
\caption{Implementation Status}
\end{table}

\clearpage

\section{Integration Workflow}

\subsection{6.1 Current Usage}

\textbf{Step 1: Setup Environment}
\begin{lstlisting}[language=bash]
# Clone repository
git clone <repository-url>
cd CD-Project

# Create virtual environment
python -m venv venv
source venv/bin/activate  # On Windows: venv\Scripts\activate

# Install dependencies
pip install -r requirements.txt

# Configure API key
echo "GEMINI_API_KEY=your_api_key_here" > .env
\end{lstlisting}

\textbf{Step 2: Test Lexer}
\begin{lstlisting}[language=bash]
# Run manual tests
python manual_test.py

# Run unit tests
pytest tests/test_lexer.py -v
\end{lstlisting}

\textbf{Step 3: Test AI Integration}
\begin{lstlisting}[language=bash]
python src/tests/ai_test_cases.py
\end{lstlisting}

\subsection{6.2 Example Output}

\textbf{Input Code (with intentional error):}
\begin{lstlisting}[language=Python]
def factorial(n)
    if n == 0:
        return 1
    return n * factorial(n - 1)
\end{lstlisting}

\textbf{Lexer Output:}
\begin{lstlisting}
Token(DEF, 'def', line=1, col=0)
Token(IDENTIFIER, 'factorial', line=1, col=4)
Token(LPAREN, '(', line=1, col=13)
Token(IDENTIFIER, 'n', line=1, col=14)
Token(RPAREN, ')', line=1, col=15)
ERROR: Missing colon ':' at line 1, column 16
\end{lstlisting}

\textbf{AI Response (from Gemini):}
\begin{quote}
\textit{"The function definition is missing a colon (:) at the end of line 1. According to the Mini-Python grammar, a function definition must follow the pattern:}

\texttt{func\_def = "def" , identifier , "(" , [ params ] , ")" , ":" , block ;}

\textit{The colon is required before the indented block begins. Corrected line should be:}

\texttt{def factorial(n):}
\end{quote}

\clearpage

\section{Project Metrics and Statistics}

\begin{table}[H]
\centering
\begin{tabularx}{\textwidth}{|l|X|}
\hline
\rowcolor{blue!10}
\multicolumn{2}{|c|}{\textbf{Overall Project Statistics}} \\
\hline
Total Weeks Planned & 16 weeks \\
\hline
Weeks Completed & 5 weeks (31.25\% complete) \\
\hline
Total Python Files & 8 files \\
\hline
Total Lines of Code & ~1,200 lines (excluding tests) \\
\hline
Test Files & 2 files \\
\hline
Test Cases & 50+ tests \\
\hline
Documentation Files & 10 files (Markdown + LaTeX) \\
\hline
\end{tabularx}
\caption{Project Statistics}
\end{table}

\begin{table}[H]
\centering
\begin{tabularx}{\textwidth}{|l|c|c|}
\hline
\rowcolor{blue!10}
\textbf{Component} & \textbf{Status} & \textbf{Completion \%} \\
\hline
Lexer & \statusgood{COMPLETE} & 100\% \\
\hline
Parser & \statusissue{NOT STARTED} & 0\% \\
\hline
Semantic Analyzer & \statusissue{NOT STARTED} & 0\% \\
\hline
AI Integration & \statuspending{PARTIAL} & 60\% \\
\hline
Web UI & \statusissue{NOT STARTED} & 0\% \\
\hline
Testing Framework & \statuspending{IN PROGRESS} & 70\% \\
\hline
Documentation & \statuspending{IN PROGRESS} & 75\% \\
\hline
\end{tabularx}
\caption{Component-wise Completion Status}
\end{table}

\section{Conclusion and Next Steps}

\subsection{8.1 Key Achievements}

This project has successfully completed the foundational phases of an Agentic AI Compiler Assistant:
\begin{itemize}
    \item \textbf{Novel architecture} combining traditional compiler design with AI assistance
    \item \textbf{Production-quality Lexer} with 98\% test coverage
    \item \textbf{Working AI integration} capable of syntax error detection and security auditing
    \item \textbf{Comprehensive documentation} of requirements, architecture, and implementation
\end{itemize}

\subsection{8.2 Immediate Next Steps (Week 6)}

\textbf{Focus:} Parser and AST Construction
\begin{enumerate}
    \item Define AST node classes for all grammar constructs
    \item Implement recursive descent parser
    \item Handle operator precedence and associativity
    \item Create comprehensive parser test suite
    \item Integrate parser with existing Lexer
\end{enumerate}

\subsection{8.3 Long-term Roadmap}

\begin{itemize}
    \item \textbf{Weeks 7-8:} Semantic analysis and symbol tables
    \item \textbf{Weeks 9-10:} Orchestrator and full AI integration
    \item \textbf{Weeks 11-12:} Web UI with Streamlit
    \item \textbf{Weeks 13-14:} Advanced features (code optimization suggestions)
    \item \textbf{Weeks 15-16:} Final testing, documentation, and deployment
\end{itemize}

\vspace{10mm}
\noindent\textbf{Document Prepared By:} Project Team \\
\textbf{Last Updated:} Week 5 Completion \\
\textbf{Document Version:} 1.0

\end{document}
